\section{Introduction}\label{sec:intro}

Fix the total wrapping charge of membranes on a curve. What accounts
for the dependence of the topological-string partition function on its
coupling at that charge? In the stable-pair calculation on a local
Calabi--Yau threefold, this dependence is encoded by a variable $q$
that grades the Euler characteristic of the sheaf: a pair of Euler
characteristic $n$ has weight $(-q)^n$. In the fivefold description
proposed by Nekrasov and Okounkov \cite{no}, the same $q$ is the
fugacity for a circle rotating the two complex directions normal to
the threefold. Selecting a different circle in the fivefold can change
the threefold on which one counts. Can the rotation weights of the new
geometry recover the Euler-characteristic series of the old one?

We answer this question for two explicit circles over a smooth curve
of positive genus. One fixes a local Calabi--Yau threefold in which the
curve has two normal line bundles of equal degree. The other fixes
the product of the curve with a plane. On the product threefold, $q$
rotates the plane, and a different variable counts the stable pairs.
We show why this second counting variable disappears and how the
weights of the plane rotation reproduce the entire $q$-dependence in
each curve degree. The two calculations agree already on the space of
curves in the fivefold, as the normal line bundle varies. This is a
case of the independence of the external circle formulated by Holmes
and Schuler \cite[Conjecture~1.7 and Examples~1.13--1.14]{holmesschuler}.

The distinction between the two variables of the topological string
will be important. We write $Q^m$ for curve degree $m$ and $q=e^{iu}$
for the exponential of the string coupling $u$. The wrapping variable $Q$ is common to both descriptions. In the Gopakumar--Vafa formula, the dependence on $u$
at fixed total degree comes from spin traces and repeated windings of
wrapped particles \cite{gv,dw}. Here we derive that dependence by
comparing two geometric indices. The doubled curve in
section~\ref{sec:collision} will let us ask how much this calculation
can tell us about a particle of charge two.

\subsection{Two threefolds in one fivefold}
Fix a smooth projective complex curve $C$ of genus $g\geq1$, put $b=g-1$,
and let $B=\operatorname{Pic}^{b}(C)$. Over $C\times B$ choose the
Poincar\'e line $\mathcal L$, trivialized along a chosen point of $C$,
and set $\mathcal M=K_C\otimes\mathcal L^{-1}$. For a line bundle $U$
and a torus character $w$, write $Uw$ for the bundle with fiber weight $w$.
Our family of fivefolds and its two fixed threefolds are
\begin{align}
 Z&=\operatorname{Tot}_{C\times B}
       (\mathcal Lt\oplus\mathcal Mt^{-1}\oplus\mathcal Oq\oplus\mathcal Oq^{-1}),
       \nonumber\\
 X=Z^{\mathbb C_q^\times}
   &=\operatorname{Tot}_{C\times B}(\mathcal L\oplus\mathcal M),
 &X'=Z^{\mathbb C_t^\times}
   &= C\times B\times S,\qquad S=\mathbb A^2 .
 \label{eq:fixed}
\end{align}
The dimensions are relative to $B$. Here a circle is a fiber rotation,
written algebraically as its complexification $\mathbb C^\times$. On $X$
the internal circle rotates the two line bundles with weights $t,t^{-1}$.
On $X'$ it rotates the plane with weights $q,q^{-1}$. Both external
circles preserve the fivefold volume form. The canonical line of $X'$ is
the pullback of $K_C$; its two normal lines in $Z$ supply the compensating
determinant.

The first fixed threefold is the balanced antidiagonal local curve used
to construct universal curve-counting classes supported on the multiples
of its zero section \cite[\S2]{wuebben-antidiagonal}. As the multiplicity
and genus vary, including genus zero, these classes freely generate the
part of Pardon's integral curve-counting algebra in first
Chern number and virtual dimension zero
\cite[Theorem~1.1]{wuebben-integral}. Embedding this model in $Z$ supplies
a second fixed threefold and a common space of cycles on which to
compare the two stable-pair indices.

A stable pair consists of a pure one-dimensional sheaf $\mathcal F$ with
proper support and a section $\mathcal O\to\mathcal F$ whose cokernel is
zero-dimensional \cite{pt}. Its curve degree $m$ is the degree of its
fundamental cycle over $C$, and $n=\chi(\mathcal F)$.
We use the modified virtual structure sheaf of \cite{no}.
The virtual structure sheaf counts with the deformation and obstruction
spaces of the pair. Its modification uses a square root of the virtual
canonical line and the determinant of cohomology from the two ordered
external lines. Exchanging the circles changes those lines as well as the
fixed threefold. We specify both roots in sections~\ref{sec:original}
and~\ref{sec:rotated}. Throughout the comparison we work relative to the
fixed curve, without integrating over $B$.

\subsection{The common index}

Curve degree one illustrates the exchange. On the product threefold a
copy of $C$ can sit at any point of the plane. At the minimum Euler
characteristic $1-g$, these translates are all the pairs. Their
deformations are the two plane translations, and their obstructions
are $g$ copies of the same two directions, one for each element of a
basis of $H^1(C,\mathcal O_C)$. The only fixed translate sits at the
origin. Cancelling the deformation factor against one obstruction
factor leaves the rotation character
\[
             Z_1(q)=\big[(1-q)(1-q^{-1})\big]^{g-1}.
\]
For a genus-two curve this is $2-q-q^{-1}$. In the original threefold
the three powers of $q$ belong to three Euler characteristics. In the
plane they all come from rotations at a single fixed pair. The
all-degree theorem extends this observation to nonreduced curves and
keeps track of their support.

Let $W_m$ be the relative Chow space of compact effective one-cycles of
degree $m$ in $Z/B$. Stable pairs on either fixed threefold determine
points of this space by taking their fundamental cycles in $Z$.
Write $s_m:B\to W_m$ for the section $m[C]$. The group
$G_0^{q,t}(W_m)$ records equivariant coherent sheaves modulo exact
sequences, with their integral multiplicities and torus characters.
Let $\mathcal I_m^{(q)}$ and $\mathcal I_m^{(t)}$ be the series of the
proper direct images of the modified sheaves in the two presentations.
Our counting convention gives weight $(-q)^n$ in the first. In the
second the weight is $(-t)^{n+mb}$, including the shift for the ordered
external line of degree $b$. The theorem in particular makes both sums
finite in this group.

\begin{theorem}\label{thm:exchange}
For the family \eqref{eq:fixed}, with the ordered determinant orientations
above, the two stable-pair indices have the same class in
$G_0^{q,t}(W_m)$ for each $m\geq 1$. More precisely,
\begin{equation}
 \mathcal I_m^{(q)}=\mathcal I_m^{(t)}
 =(s_m)_*\bigl(Z_m(q)[\mathcal O_B]\bigr),\qquad
 Z_m(q)=\sum_{\lambda\vdash m}
      \prod_{\square\in\lambda}
        \big[(1-q^{h(\square)})(1-q^{-h(\square)})\big]^{g-1}.
 \label{eq:hook}
\end{equation}
Here $h(\square)$ is one plus the number of boxes to the right and below
$\square$ in its row and column. The equality holds over the full Picard
variety, including the locus where the number of sections of the normal
bundle jumps.
\end{theorem}

On a fixed target the character $Z_m(q)$ is known. Bryan and
Pandharipande obtained the local Gromov--Witten partition function
\cite{bryanpandharipande}, and Monavari's stable-pair formulas give
its $q$-character \cite{monavari,refined}. For the related comparison
with fivefold Gromov--Witten theory, Giacchetto, Pandharipande and
Schuler derive a degree-one local-curve correspondence from their
membrane conjectures \cite{gps}. The content of
Theorem~\ref{thm:exchange} is that two moduli problems give the
same integral class supported on $m[C]$. Equality of the Laurent
polynomials on individual targets would leave this undetermined:
a nonzero class on the Picard variety can have rank zero at every
point.

The two sides account for the equality in different ways. On $X$
the infinitesimal rotation of the normal lines gives a cosection of
the obstruction sheaf. It concentrates the virtual sheaf on invariant
pairs. Their fixed deformation theory, with its determinant orientation,
is constant over $B$, even where the dimension of the space of normal
sections jumps (section~\ref{sec:original}).

On $X'$ the two plane coordinates act as endomorphisms of a rank-$m$
bundle $E$ on $C$. Its degree $k=n+m(g-1)$ is precisely the exponent
of the exchanged counting variable. Write $x,y$ for the endomorphisms.
The framed matrix equation admits the deformation
\[
                      [x,y]+ij=a\,1_E ,
\]
where $i$ is the section of the pair, $j$ a covector, and $a$ a scalar.
At $a\ne0$, taking traces on invariant quotients and subspaces forces
the words in $x,y$ applied to $i$ and to $j$ to generate $E$ and
$E^\vee$, respectively. Both bundles are therefore generated by global
sections. Their degrees must both be nonnegative, so $k=0$.
For $k>0$, specialization back to the original equation makes the
virtual structure sheaf vanish on the stable-pair space itself. At
$k=0$ the matrix data are constant along $C$ and specify a length-$m$
subscheme of the plane. The tangent weights of its Hilbert scheme
supply the hook lengths in \eqref{eq:hook}
(section~\ref{sec:rotated}). This accounts for the loss of the
$t$-grading and the survival of the $q$-dependence in the same calculation.

\subsection{What the doubled curve detects}
In genus two, degree two is the first place where two independent
degree-one objects contribute at the same total charge as a connected
double. Over two distinct points of the plane the index sheaf factors.
This prescribes the graded symmetric square to subtract for the
independent pair. The remaining class is supported where the two
points coincide. There the rotation has two fixed length-two schemes,
one for each coordinate direction in which a point can thicken. Each
contributes $+1$ at $q^{-3}$ and $-2$ at $q^{-2}$. The graded symmetric
square contributes zero at both powers, giving the connected
coefficients $2$ and $-4$ (section~\ref{sec:collision}).
To interpret the latter as the spin trace of a charge-two particle,
one would also have to establish a normalizable state in the wrapped
theory with the continuum of two separated branes open and derive its
spin.

Varying the product $\kappa$ of the normal weights on $X$ detects
the principal polarization of the Picard variety. Give the four lines
in \eqref{eq:fixed} the weights $t,\kappa/t,q,\kappa^{-1}/q$, so their
product remains one. Write $\delta=1-\kappa$.

\begin{theorem}[First variation in genus two and curve degree two]\label{thm:picard-variation}
Let $C$ have genus two, let $B=\operatorname{Pic}^1(C)$, and put
$A=R\pi_*\mathcal L$ for the normalized Poincar\'e line on
$C\times B$, with $\pi:C\times B\to B$. With the ordered orientations
of Theorem~\ref{thm:exchange}, the original and rotated family indices
have the same coefficient of $\delta q^{-3}t^k$ for every $k\geq1$.
Its value in $K^0(B)$ is
\[
       -2\bigl(1+2\,\mathbf1_{2\mid k}\bigr)\psi^k(A),
\]
where $\psi^k$ is the $k$th Adams operation. On $X$ one takes the
$q^{-3}$ coefficient and expands its first variation at $t=0$.
On $X'$ one takes the $t^k$ coefficient and expands its first
variation at $q=0$.
\end{theorem}

At $k=2$ the Chern character of this class is $12\theta$, where
$\theta$ is the principal polarization of $B$. On $X$ it comes from
two ribbons, one in each normal line. On $X'$ the two types of divisor
contribute $4\theta$ and $8\theta$. Restriction to any individual
target gives zero because $\chi(C,L)=0$ for every $L\in B$.
Section~\ref{sec:variation} gives both calculations at the lowest
original grading $q^{-3}$.
